%% ma: L10-B14-0036
%% lop: 10
%% chuong: 5
%% bai: 14
%% dang: 10.14.1
%% loai: TLN
%% muc: NB
%% dap_an: "13"
%% nguon: "chuyen-de-cac-so-dac-trung-cua-mau-so-lieu-khong-ghep-nhom-toan-10-chuong-trinh-moi.pdf (GV. Phan Nhật Linh) (P04), Bài 03 Dạng 1 Phần III Câu 3"
%% kien_thuc: "Lớp 10 Bài 14 (khoảng biến thiên)"
%% trang_thai: chinh-thuc
%% ghi_chu: "Tự duyệt theo yêu cầu giáo viên 22/10/2026."
\begin{ex}
Mẫu số liệu sau cho biết chiều cao (đơn vị: cm) của các bạn trong tổ: $163$; $159$; $172$; $167$; $165$; $168$; $170$; $161$. Tìm khoảng biến thiên của mẫu số liệu.
\shortans{13}
\loigiai{$R=172-159=13$.}
\end{ex}
